\chapter{The Big Rupture: The Mechanical Origin of the Observable Universe}
\label{chap:the_big_rupture}

Standard cosmology posits that the universe began from a singularity---a mathematical point of infinite density and temperature---expanding outward in a ``Big Bang.'' Tension Field Theory (TFT) resolves this non-physical paradox by replacing the singularity with a strictly mechanical crystallographic phase transition.

The genesis event was not an explosion of ``nothing'' into ``something.'' It was the Catastrophic Decompression of a saturated region of the spacetime lattice.

\section{The Three States of the Spacetime Lattice}
To understand the mechanics of the Big Rupture, we must recall the foundational structural phases of the Spacetime Gravitational Floor (SGF) established in Part I. The universe operates strictly through phase transitions between three mechanical tension states ($\tau$):

\begin{itemize}
    \item \textbf{State 0: The Latent Substrate (Equilibrium, $\tau = \tau_S$).} The absolute, invariant, and infinite background frame of reality. It is a perfectly rigid, Time-Frozen Simple Cubic (SC) crystal with zero interstitial voids and zero entropy ($S = 0$).
    \item \textbf{State I: The Active Vacuum (Expansion, $\tau = \tau_D$).} The dynamic, time-enabled ``Spacetime'' we observe. It acts as an elastic skin over the rigid substrate, configured in a Face-Centered Cubic (FCC) geometry. Its interstitial voids provide the kinematic freedom necessary for transverse waves (light) and hydrodynamic inertia (mass).
    \item \textbf{State II: The Saturated Core (Compression, $\alpha_G \to 1$).} A localized, macroscopic collapse of the State I active vacuum into the State II saturated phase. When gravitational tension reaches absolute saturation ($\alpha_G = 1$), the active FCC voids are crushed out of existence, mechanically arresting time. The region topologically locks into the hyper-dense, static SC configuration.
\end{itemize}

The Observable Universe is exclusively a State I phenomenon. The Big Rupture is the mechanical mechanism by which a global network of State II cores violently decompresses. Because the Black Hole previously consumed the surrounding State I vacuum, its catastrophic decompression does not expand into ``nothingness.'' Instead, it snaps outward into the exposed, invariant State 0 background, mechanically forcing interstitial voids back into the latent structure to restore the State I active vacuum.

\section{The Pre-Rupture Trigger (Global Topological Percolation)}

As derived in Part II, a Black Hole (State II core) grows by consuming the active State I vacuum, converting FCC nodes into the void-less SC Substrate ($\rho_S \approx 10^{96} \text{ kg/m}^3$). Any pre-existing dynamic mass-energy is mechanically displaced outward, forming a hyper-compressed FCC boundary shell (the Event Horizon) that is stabilized strictly by the inward Boundary Tension ($\tau_{QG}$) and supported by the ambient pressure of the surrounding active vacuum.

However, a Black Hole does not grow in isolation. The universe contains a distributed, expanding network of these macroscopic State II defect cores. While a single theoretical core possesses a local causal growth limit ($M_{\text{Sat}}$), the universe never reaches this localized extreme. The rupture is strictly triggered by global geometric saturation: \textbf{Site Percolation}.

As the macroscopic network of frozen State II cores accumulates, it consumes an increasing percentage of the active vacuum's interstitial capacity. When this global defect network reaches exactly 5.16\% of the reservoir's capacity (detailed in \hyperref[app:energy_and_chronology]{Appendix~\ref*{app:energy_and_chronology}}), the isolated cores mechanically connect to form a contiguous, macroscopic fault line across the entire cosmos.

This global topological fracture physically shatters the structural continuity of the active State I vacuum. Consequently, the ambient containment pressure universally vanishes ($\Delta P \to 0$). Without this external pressure to support the Event Horizons, the hyper-compressed boundary shells simultaneously mechanically dissolve.

\section{Trigger Boundary Parameters at Percolation}
\label{sec:trigger_boundary_parameters}

The Catastrophic Decompression is governed by a strict set of global boundary parameters that define the state of the macroscopic reservoir at the exact moment the 5.16\% percolation limit is breached. As the active Face-Centered Cubic (FCC) vacuum is consumed and converted into solid Simple Cubic (SC) cores, the macroscopic boundary must geometrically deflate to maintain contiguous tension.

As rigorously derived and geometrically proven in \hyperref[app:energy_and_chronology]{Appendix~\ref*{app:energy_and_chronology}}, the reservoir reaches absolute structural failure under the following global limits:

\begin{itemize}
    \item \textbf{Global Rupture Trigger Mass ($M_{\text{trg}}$):} The rupture is strictly triggered when the Aggregate Percolating Mass ($M_p$)—the total mass-energy segregated onto the global network of State II phase boundaries—reaches exactly $5.16\%$ of the total macroscopic reservoir mass ($\approx 1.02 \times 10^{78} \text{ kg}$).
    \item \textbf{Deflated Rupture Radius ($R_{\text{rup}}$):} The physical boundary of the Geometric Cosmic Horizon deflates by approximately $1.09 \times 10^{33} \text{ m}$, locking the final rupture radius at $\approx 6.077 \times 10^{34} \text{ m}$.
    \item \textbf{Global Pressure Drop ($\Delta P \to 0$):} At this exact density, the distributed SC cores physically intersect shoulder-to-shoulder, forging a contiguous, unbroken fractal fault line across the entire macroscopic diameter of the system. This rigid macroscopic skeleton structurally shatters the elastic continuity of the remaining 94.84\% FCC lattice, dropping ambient containment pressure to zero and triggering the universe-wide martensitic phase transition.
\end{itemize}

\begin{figure}[H]
    \centering
    \includegraphics[width=0.9\textwidth]{Fig20_PartIII_BigRupture.pdf}
    \caption{\textbf{Catastrophic Decompression via Global Percolation (The Big Rupture).} \textbf{(a) At Rupture:} The macroscopic cosmic reservoir deflates ($R_{\text{rup}} \approx 6.077 \times 10^{34} \text{ m}$) as active vacuum is consumed by State II SC Cores (Black Holes). At precisely $5.16\%$ volumetric capacity, the isolated cores geometrically percolate. This forms a contiguous, rigid spanning cluster (Global Percolation Fault Line) that universally shatters the containment pressure ($\Delta P \to 0$). \textbf{(b) Post Rupture:} Stripped of boundary tension, the hyper-compressed SC fault line violently decompresses, expanding outward to restore the active FCC vacuum geometry ($R_{\text{Res}} \approx 10^{35} \text{ m}$). The phase-segregated Baryonic Matter (Topological Defects) previously trapped on the Event Horizons is explosively expelled, surfing the expanding longitudinal shockwaves as the Phonon Echo (BAO rings) propagates through the restored lattice.}
    \label{fig:big_rupture}
\end{figure}

\section{The Big Rupture (Catastrophic Decompression)}

With the global network of Event Horizons destroyed, the contiguous fault line of hyper-compressed State II cores is suddenly exposed without a containment gradient. Because the Simple Cubic packing factor (1.0) is a state of extreme mechanical stress compared to the relaxed Face-Centered Cubic packing factor ($\approx 0.74$), the contiguous defect network violently yields.

The entire macroscopic State II fault line undergoes a martensitic phase transition. We define this sudden, diffusionless structural snap-back of the lattice as Catastrophic Decompression. The zero-void SC lattice explosively snaps back into the actively spaced FCC geometry. This instantaneous, volumetric expansion of the lattice phase from State II back into State I is the mechanical reality of the Big Bang.

It is not an expansion of empty space; it is the physical snap-back of a hyper-compressed crystal restoring its interstitial voids. Because this martensitic (diffusionless) phase transition propagates via the longitudinal stress transmission of the underlying substrate ($c_S \to \infty$), it is not bound by the transverse shear wave limit of light ($c_E$). Therefore, the entire macroscopic fault line snaps back simultaneously. This longitudinal shockwave effortlessly replaces the need for an abstract ``inflation field'' by mechanically guaranteeing a uniform temperature and density distribution across the macroscopic reservoir from the very first moment.

\enlargethispage{\baselineskip}
The universe we inhabit today was born in the fraction of a second when the global topological fault line shattered the vacuum, forcing the timeless structural sinks to violently rupture back into the active flow of time.
